\documentclass[10pt,a4paper]{amsart}
\usepackage[margin=19mm]{geometry}
\usepackage{amsmath,amssymb,mathtools}
\usepackage[scr=rsfs,cal=euler]{mathalfa}
\usepackage{xfrac}
\usepackage[unicode,hidelinks,bookmarks=false]{hyperref}
\newcommand{\ie}{i.e.\ }
\newcommand{\cf}{cf.\ }
\newcommand{\resp}{respectively\ }
\newcommand{\Subsection}{\subsection}
\providecommand{\nobreakdash}{\nobreak\hbox{-}}
\setlength{\emergencystretch}{2em}
\multlinegap0pt
\usepackage{bm}
\usepackage{bbm}
\usepackage{dsfont}
\usepackage{xspace}
\usepackage{cancel}
\usepackage{mathtools}
\usepackage[shortlabels]{enumitem}
\usepackage{amsthm}
\makeatletter
\@ifundefined{thm}{\newtheorem{thm}{Thm}}{}
\@ifundefined{prop}{\newtheorem{prop}[thm]{Proposition}}{}
\makeatother
\makeatletter
\def\NN{\mathbb N}
\expandafter\ifx\csname Proof\endcsname\relax
\newenvironment{Proof}[1]{ \raggedright \textbf{Proof of {#1}.} }{\qedwhite}
\fi
\def\RR{\mathbb R}
\def\ZZ{\mathbb Z}
\def\esssup{\mathop{\textup{ess}\sup}}
\expandafter\ifx\csname proof\endcsname\relax
\renewenvironment{proof}{{\raggedright \bfseries Proof.}}{\qedwhite}
\fi
\newcommand{\vc}[1]{\overrightarrow{#1}}
\makeatother
\title[Constructive Approximation in Mixed norm Spaces]{Constructive Approximation in Mixed norm Spaces}
\author{Priyanka Majethiya, Shivam Bajpeyi, Dhiraj Patel}
\date{Source: arXiv:2506.04787v2 (CC BY 4.0)}
\begin{document}
\maketitle

\noindent\emph{Source and licence.} Constructive Approximation in Mixed norm Spaces. Version arXiv:2506.04787v2 (CC BY 4.0), distributed under the Creative Commons Attribution 4.0 International licence, as recorded in the arXiv licence metadata for this exact version and shown on the abstract page https://arxiv.org/abs/2506.04787v2. Full rights notice: licenses/arxiv-2506.04787-CC-BY-4.0.txt.\par\smallskip

\noindent\textit{Editorial scope.} Counted unit: the Prop whose source label is \texttt{p1} in the article (source0.tex lines 386--388, source label \texttt{p1}), delivered together with the article's own complete original proof of it (source0.tex lines 389--427, one \texttt{proof} environment of 39 lines). No mathematics is written, repaired or re-derived: statement and proof are the article's own text line for line. Required context retained uncounted: none. Every definition, display and label of those lines is retained unchanged. Omitted from this package and not redistributed: 1--385, 428--1295. Nothing inside a retained excerpt is omitted or altered: every retained body is delivered exactly as the article's lines are written, so no omission receipt of any kind accompanies this package. Citations: the retained excerpts carry no citation command at all, so no bibliography entry of the article is transcribed and no reference is redistributed. Reference adaptation: none was needed and none was made; every reference inside the delivered unit resolves inside it, so no unresolved reference or citation is produced. Printed slips kept literal: no printed word, symbol, hypothesis or number of the counted statement or of its proof was altered. Layout and class: the article's own class is \texttt{article} with packages including amsmath, amsthm, amscd, amsfonts, amssymb, enumitem, latexsym, multirow, xcolor, lipsum, comment, hyperref; the wrapper is the lane's pinned \texttt{amsart} editorial wrapper at 10pt on A4, which restates the theorem-like environments the retained text needs and adds the article's own macro definitions that the retained text needs (\texttt{\textbackslash NN}, \texttt{\textbackslash RR}, \texttt{\textbackslash ZZ}, \texttt{\textbackslash esssup}, \texttt{\textbackslash vc}), with extra wrapper packages (bm, bbm, dsfont, xspace, cancel, mathtools, enumitem). The delivered numbering is the unit's own. Assets: no retained span contains a figure environment, a graphics-inclusion command, a PDF-page inclusion, a table or a TikZ picture; no image file is copied into this package, the package declares no asset dependency and no image file is redistributed with it. Licence: version arXiv:2506.04787v2 of this article is distributed under the Creative Commons Attribution 4.0 International licence, as recorded in the arXiv licence metadata for this exact version and shown on the abstract page https://arxiv.org/abs/2506.04787v2. A full rights notice is delivered at licenses/arxiv-2506.04787-CC-BY-4.0.txt. Characters: every retained excerpt is pure ASCII, so the delivered engine, XeLaTeX with its default Latin Modern text font, reports no missing character.
\begin{prop} \label{p1}
    The space $\Delta^{\vc{P}}(\RR^n)$ is a proper subspace of $L^{\overrightarrow{P}}(\RR^n)$, and $\|\cdot\|_{\Delta^{\vc{P}}}$ is a seminorm on $\Delta^{\vc{P}}(\RR^n)$.
\end{prop}

\begin{proof}
    We begin by demonstrate the result for $n=2$ and then generalize it for any $n \in \NN$.
    We only need to show that $f \in \Delta^{\overrightarrow{P}}$ implies $\displaystyle \left(\int_{\RR}\left(\int_{\RR} |f(x_1,x_2)|^{p_1}dx_1 \right)^{\frac{p_2}{p_1}}dx_2 \right)^{\frac{1}{p_2}}$ is finite. Since $f$ is bounded, it is possible to find $x_{k_1} ^{*}\in [k_1,k_1+1],y_{k_2}^{*}\in [k_2,k_2+1]$, for each $k_1,k_2 \in \ZZ$, such that 
    \[ \esssup_{x_1 \in [k_1,k_1+1], \, x_2 \in [k_2,k_2+1]} |f(x_1, x_2)| < |f(x_{k_1}^{*},y_{k_2}^{*})| + \frac{1}{1+k_1^2+k_2^2},\]
    and it follows that
    \begin{align*}
        \|f\|_{(p_1,p_2)} & =\Bigg(\sum_{k_2 \in\ZZ}\int_{k_2}^{k_2+1}\Bigg(\sum_{k_1 \in \ZZ} \int_{k_1}^{k_1+1}|f(x_1,x_2)|^{p_1}dx_1\Bigg)^{\frac{p_2}{p_1}}dx_2\Bigg)^{\frac{1}{p_2}}\\
		& < \Bigg( \sum_{k_2 \in \ZZ}\Bigg(\sum_{k_1 \in \ZZ}\Bigg(|f(x_{k_1}^{*},y_{k_2}^{*})|+ \frac{1}{1+k_1^2+k_2^2}\Bigg)^{p_1}\Bigg)^{\frac{p_2}{p_1}}\Bigg)^{\frac{1}{p_2}}\\
		& \leq \Bigg(\sum_{k_2 \in \ZZ}\Bigg( \sum_{k_1 \in \ZZ}|f(x_{k_1}^{*},y_{k_2}^{*})|^{p_1}\Bigg)^{\frac{p_2}{p_1}}\Bigg)^{\frac{1}{p_2}}+ \Bigg(\sum_{k_2 \in \ZZ}\Bigg(\sum_{k_1 \in \ZZ} \Bigg( \frac{1}{1+k_1^2+k_2^2}\Bigg)^{p_1}\Bigg)^{\frac{p_2}{p_1}}\Bigg)^{\frac{1}{p_2}}  \\
		&:= S_1+S_2 \hspace{2pt}.
    \end{align*}
    Indeed, $S_2$ is finite for each $p_1,p_2\geq 1$. Concerning $S_1$, the problem is that $(x_{k_1}^{*},y_{k_2}^{*})_ {(k_1,k_2) \in \ZZ^2}$ may not necessarily be a relatively separated in $\RR^2$, but the subsequences $(x_{2k_1}^{*},y_{2k_2}^{*})_ {(k_1,k_2) \in \ZZ^2}$, $(x_{2k_1+1}^{*},y_{2k_2}^{*})_ {(k_1,k_2) \in \ZZ^2}$, $(x_{2k_1}^{*},y_{2k_2+1}^{*})_ {(k_1,k_2) \in \ZZ^2}$, $(x_{2k_1+1}^{*},y_{2k_2+1}^{*})_ {(k_1,k_2) \in \ZZ^2}$ are relatively separated set in $\RR^2$, and hence for each $f \in \Delta^{\overrightarrow{P}}$, we have
    \begin{align*} 
        &\Bigg( \sum_{k_2 \in \mathbb{Z}}\Bigg( \sum_{k_1 \in \mathbb{Z}} |f(x_{k_1}^*, y_{k_2}^*)|^{p_1} \Bigg)^{\frac{p_2}{p_1}}\Bigg)^{\frac{1}{p_2}}\\
        \leq& \Bigg( \sum_{k_2 \in \mathbb{Z}}\Bigg(\sum_{k_1 \in \mathbb{Z}}  |f(x_{2k_1}^*, y_{2k_2}^*)|^{p_1} \Bigg)^{\frac{p_2}{p_1}}\Bigg)^{\frac{1}{p_2}} + \Bigg(\sum_{k_2 \in \mathbb{Z}}\Bigg(\sum_{k_1 \in \mathbb{Z}}   |f(x_{2k_1}^*, y_{2k_2+1}^*)|^{p_1} \Bigg)^{\frac{p_2}{p_1}}\Bigg)^{\frac{1}{p_2}} \\
        &\quad + \Bigg(\sum_{k_2 \in \mathbb{Z}} \Bigg( \sum_{k_1 \in \mathbb{Z}} |f(x_{2k_1+1}^*, y_{2k_2}^*)|^{p_1} \Bigg)^{\frac{p_2}{p_1}}\Bigg)^{\frac{1}{p_2}} + \Bigg( \sum_{k_2 \in \mathbb{Z}}\Bigg(\sum_{k_1 \in \mathbb{Z}}  |f(x_{2k_1+1}^*, y_{2k_2+1}^*)|^{p_1} \Bigg)^{\frac{p_2}{p_1}}\Bigg)^{\frac{1}{p_2}}\\
        <&\, \infty.
    \end{align*}
    In a similar manner, for $n\in \NN$ there exist $2^n$ subsequent relatively separated sets. Using the same reasoning, since  $f\in \Delta^{\vc{P}}(\RR^n)$, it follows that the all the mixed norm Lebesgue space are finite. Hence, $f\in L^{\vc{P}}(\RR^n)$. 
    For the second part, we provide the proof in the following steps:
    \begin{enumerate}[label=(\roman*), left=0pt]
        % \item [(i)] According to the definition of $\|f(x_{\textbf{k}})\|_{\ell^{\overrightarrow{P}}_w}$, it is clear that it is non-negative. Consider \[ f(\textbf{u}) = 
        % \begin{cases}
        %     0, & \text{if } \textbf{u} = \textbf{k} \text{ for } \textbf{k} \in \mathbb{Z}^n,\\
        %     1, &  otherwise,
        % \end{cases}\]
        % we have  $\|f(x_{\textbf{k}})\|_{\ell^{\overrightarrow{P}}_w}=0,$ but $f$ is non-zero.

        \item For any $f\in \Delta^{\vc{P}}(\RR^n)$, we have 
       % $ p(\alpha f) = |\alpha| \hspace{2pt} p(f)$
        $\|\alpha f\|_{\Delta^{\vc{P}}}= |\alpha|  \hspace{2pt} \|f\|_{\Delta^{\vc{P}}}$,  
        $ \forall \alpha \in \RR$. 
        
        \item For any $f, g\in \Delta^{\vc{P}}(\RR^n)$ and utilizing triangle inequality, we have 
        % $$p(f+g) \leq p(f) + p(g).$$
        $$\|f+g\|_{\Delta^{\vc{P}}} \leq \|f\|_{\Delta^{\vc{P}}}+\|g\|_{\Delta^{\vc{P}}}.$$
    \end{enumerate}
    Therefore, $\|\cdot\|_{\Delta^{\vc{P}}} $ is a seminorm on $\Delta^{\vc{P}}(\RR^n)$. This completes the proof.
\end{proof}

\end{document}
